\documentclass[10pt,a4paper]{article}
\usepackage[utf8]{inputenc}
\usepackage[frenchb]{babel}
\usepackage[T1]{fontenc}
\usepackage{amsmath}
\usepackage{amsfonts}
\usepackage{amssymb}
\usepackage{graphicx}
\usepackage{lmodern}
\usepackage[left=2cm,right=2cm,top=2cm,bottom=2cm]{geometry}
\author{Alain Fontan}
\title{Chronologie du projet}
\begin{document}
\section{Cheminement de l'élaboration du questionnaire}

\subsection{Élaboration d'un projet de questionnaire par itération entre les membres du Comité pour une version stabilisée en juin 2026}

\textbf{Réunions et initiatives du Comité :}

\begin{itemize}
\item 21/11/2025
\begin{itemize}
\item Identification d'axes de travail sur la place des réserves citoyennes et stratégiques,
\item Rôle des compétences issues de l'IHEDN dans la résilience locale, la circulation de l'information, la sensibilisation de la population,
\item Les interactions entre associations régionales et autorités publiques,
\item Le besoin de coordination, de formation et de partage d'expérience,
\item Le renforcement d'une culture territoriale de sécurité globale.
\end{itemize}
\item 11/12/2025
\begin{itemize}
\item Notre réflexion s'oriente désormais plus directement vers le rôle opérationnel que peuvent jouer les citoyens engagés et les associations régionales de l'IHEDN dans la prévention, la gestion et la résilience face aux crises majeures.
\item Proposition d'une cartographie des compétences, réflexion sur une éventuelle réserve IHEDN, renforcement des actions de sensibilisation et de la coordination des autorités locales.
\item Points de vigilance et défis à relever pour assurer la faisabilité et l'acceptabilité de ces orientations.
\end{itemize}
\item 22/01/2026
\begin{itemize}
\item Problématique du comité :  les associations régionales de l'IHEDN n'ont pas pour vocation à décrire ou gérer directement une crise, mais à contribuer utilement à la préparation, à l'accompagnement des acteurs et à l'analyse des crises, en complément des autorités compétentes.
\item L'annuaire de l'Union IHEDN représente un vivier significatif de compétences et d'expertise, susceptible de constituer le socle d'un dispositif structuré de préparation aux crises et de retour d'expérience.
\item La réflexion menée permet d'aboutir à une proposition concrète de la part de notre comité : l'élaboration d'un questionnaire à destination des Auditeurs de l'IHEDN. Ce questionnaire a pour objectif de cartographier les ressources disponibles, leurs domaines d'expertise, les profils volontaires pour intégrer un dispositif d'appui.
\item Un calendrier prévisionnel de cette démarche est envisagé après consultation de l'IRSEM.
\item Le Comité sollicitera l'accord de la Gouvernance de l'AR 16.
\end{itemize}
\item 12/02/2026
\begin{itemize}
\item Confirmation de la proposition du Comité à mieux connaître les membres de l'Association par l'intermédiaire d'une enquête rigoureuse, crédible et structurée. Le dispositif envisagé a pour objectif de mieux connaître les membres de l'AR 16, clarifier leur positionnement avant et après les crises, à proposer une intégration crédible dans les dispositifs existants et à renforcer la cohésion et la visibilité de l'Association.
\end{itemize}
\item 03/03/2026
\begin{itemize}
\item Rédaction et présentation de la problématique et du cadre général de l'enquête et de l'utilité du questionnaire.
\end{itemize}
\item 05/03/2026
\begin{itemize}
\item Formulation d'une interrogation centrale : quelle est la capacité réelle de mobilisation de l'AR 16 en cas de crise majeure ?
\item Le choix d'un questionnaire permettrait de mieux connaître les profils, compétences et disponibilité des auditeurs. Il permettrait de renforcer les liens entre les membres en manifestant l'attention de l'Association sur leurs parcours et expertises.
\item Le Cinquantième anniversaire de l'Association (mais aussi celui de l'Union-IHEDN) et les 90 ans de l'IHEDN constituent une opportunité symbolique forte pour engager une réflexion collective : l'enquête permettrait à la fois de célébrer l'histoire de l'Association et d'interroger sa capacité d'adaptation et d'engagement face aux défis contemporains.
\end{itemize}
\item 06/03/2026
\begin{itemize}
\item Dans cette perspective, une lettre de présentation générale du projet est adressée à la Présidence de l'Association AR 16 
\end{itemize}
\item 19/03/2026
\begin{itemize}
\item Suite à un échange avec le bureau directeur de l'AR 16, la mise en place d'un questionnaire semble être une idée d'intérêt pour ses membres car il « rejoint une piste évoquée sur la feuille de route AR 16… ».
\item Les échanges traduisent une convergence vers un objectif commun : mieux structurer la contribution des auditeurs des associations en conciliant réalisme (disponibilité, compétences, cadre d'action) et ambition (mobilisation et engagement opérationnel)
\item Expression d'un souhait de prolongement des travaux du comité sur une période nous menant en 2027
\end{itemize}
\item 06/03/2026
\begin{itemize}
\item Echange sur la réalité et l'appartenance à la communauté des auditrices et auditeurs. De quoi parle-t-on ? Comment s'insérer dans la multiplicité des crises ?
\item Evocation d'un entretien informel avec la Présidente de l'Union IHEDN, favorable à un questionnaire utile pour l'ensemble des Associations. A' tester à Paris-Ile-de-France, dans une région rurale (expl. ex région Limousin) et en Outre-mer.
\end{itemize}
\item 22/04/2026
\begin{itemize}
\item Communication d'une ébauche du rapport du Comité.
\end{itemize}
\item 05/05/2026
\begin{itemize}
\item Présentation aux membres du Comité de la première version du rapport du Comité « Crises majeures : quel rôle pour les citoyens engagés et les associations IHEDN ? ».
\end{itemize}
\item 19/05/2026
\begin{itemize}
\item Partage d'une analyse de l'annuaire de l'Union-IHEDN 2022 consacrée au vivier des auditeurs et aux limites du potentiel réellement mobilisable en cas de crise majeure.
\item Le rapport a été approfondi et le positionnement de l'AR 16 y est désormais davantage clarifié comme acteur de contribution indirecte mais stratégique, principalement en amont et en aval des crises, autour des enjeux de préparation, de culture stratégique, d'analyse et de retour d'expérience.
\end{itemize}
\item 21/05/2026
\begin{itemize}
\item Echange autour de l'annonce de la fin du projet d'enquête
\end{itemize}
\item 11/06/2026
\begin{itemize}
\item Relecture et derniers arbitrages sur le contenu du Rapport du Comité avant son envoi au Comité Directeur de l'AR 16.
\end{itemize}
\end{itemize}
\subsection{Identification des obstacles à surmonter}
\begin{itemize}
\item Cadre réglementaire : RGPD (applicable depuis le 25 mai 2018), encadre le traitement des données personnelles sur le territoire de l'UE et LIL (Loi Informatique et Libertés du 6 janvier 1978 modifiée), complémentaire à RGPD et mise en œuvre par la CNIL, autorité de contrôle indépendante.
\item Envoi du questionnaire numérique, de sa présentation et justification avant les étapes essentielles de relance et collecte, 
\item Traitement et analyse des données (assistance institutionnelle ou universitaire).
\item Restitution : organisation événementielle à prévoir.
\item Calendrier prévisionnel (juin 2026 -- septembre 2027)
\end{itemize}
\subsection{Recherche dans les archives/données similaires permettant de construire le projet de l'enquête et fabriquer le questionnaire}

\subsubsection{Consultation d'études nationales et régionales}

\begin{itemize}
\item Enquête nationale : Consultation du questionnaire de l'enquête nationale auprès des Associations régionales de l'Union IHEDN (Septembre 2015) dans la double perspective de :

\begin{itemize}
\item « Créer des synergies entre ses composantes, synergies qui aboutissent à ce que l'efficacité de l'Union soit supérieure à la simple addition de l'efficacité de ses membres et à ce que chaque association soit plus forte parce que liée aux autres »
\item « Renforcer globalement la communauté des auditeurs sans méconnaître les spécificités des uns et des autres ».
\end{itemize}

\item IEP de Toulouse/AR IHEDN Toulouse et Montpellier -- Rapport du C2SD* - « L'esprit de défense au quotidien : les potentialités de développement de l'IHEDN » - Enquête (anonyme) auprès des acteurs locaux (environ 300) -- Auteurs : Jacques Aben, Marie-Dominique Charlier-Dagras et Jean-Pierre Marichy (2000).
\end{itemize}

\subsubsection{Consultation de ressources bibliographiques ou en ligne}

\begin{itemize}
\item Union IHEDN : Consultation des archives ( collection  de la Revue Défense).
\item Editions Italiques : « Sécurité, Défense, société civile : nouvelles solidarités » (2004).
\item Vittorio Hösle : « Les forces centrifuges planétaires. Une cartographie du temps présent à la lumière de la philosophie de l'histoire » (2021). Traduction Bruno Leray/Christine Griesmar
\item Thèse d'histoire - « L'IHEDN : une vision globale de la politique de Défense de la France ». J. Sauvage sous la direction de Maurice Waïsse -- Reims. (1991).
\item Thèse de l'Institut Polytechnique de Paris « Intégrer les contributions citoyennes aux dispositifs de gestion de crise : l'apport des médias sociaux »   -- Robin Batard (25/06/21).
\item Thèse de doctorat de Sciences économiques et sciences de gestion : « Préparer les citoyens à la gestion des risques majeurs/risques de catastrophes à l'échelle des communes : une approche par l'instrumentation gestionnaire ». Rose Toki IAE Nantes (13/12/2021).
\item Colloque IHEDN au Sénat : « Résilience d'une nation : y-a-t-il une place pour le citoyen dans la sécurité de notre pays ? » --  Lavauzelle. Collection Les grands colloques de l'IHEDN (2018).
\item Semaine IHEDN Jeunes : AR Poitou-Charentes/Institut poitevin de recherche en sociologie de la connaissance : « La personne citoyenne » - Poitiers (31/08 -- 05/09).
\item Place Joffre : Cinquantenaire de l'IHEDN et des associations régionales d'auditeurs/IHEDN. Editions Hervas (1999).
\item Fondation pour les études de défense/Documentation française - « La gestion de sortie de crise : actions civilo-militaires et opérations de reconstruction (1998)
\item L'album des sessions : photos et liste des auditeurs des sessions nationales, régionales, européennes, africaines et malgaches (1948-1955). Editions Hervas.
\item IEP de Toulouse/AR IHEDN Toulouse et Montpellier -- Rapport du C2SD : « L'esprit de défense au quotidien : les potentialités de développement de l'IHEDN » - Enquête (anonyme) auprès des acteurs locaux (environ 300) -- Auteurs : Jacques Aben, Marie-Dominique Charlier-Dagras et Jean-Pierre Marichy (2000).
\item Les Jeunes IHEDN -- GT OPS Jeux Olympiques et Paralympiques de Paris 2024 (2023)
\item Sciences Po -- CrisisLab : Working Paper “ L'expérimentation générative et adaptative. Une approche de la préparation à la gestion de crise » (2025).
\item Centre d'analyse comparative des systèmes politiques de l'Université Paris I. Article : « les potentialités du développement des activités de l'IHEDN, enquête auprès des acteurs locaux » (1972).
\item CNRS Info -PEPR IRMa -- Article : “Il y a un besoin crucial d'innover en termes de politiques de gestion des risques et des crises », interview de Soraya Boudia (04/2024).
\item IRSEM -- Etude 70 : « Risques géopolitiques, crises et ressources naturelles -- Approches transversales et apport des sciences humaines » - Sarah Adjel, Angélique Palle, Noémie Rebière (09/2019).
\item IRSEM -- Etude 95 : « L'armée, les Français et la crise sanitaire. Une enquête inédite » (06/2022). 
\item CNRS -- Lettre Sciences Humaines et Sociales : Retour sur le séminaire « Les SHS à l'épreuve des risques et des catastrophes » (28/04/2026)
\item IHEMI - Lettre d'information sur les Risques et les Crises (LIREC) : « Comment développer la culture de crise » ? -- Anne Muxel, Florian Opillard, Angélique Palle (12/2023).
\item IHEMI -- Lettre d'information sur les Risques et les Crises (LIREC) : « Le volontariat et la mobilisation citoyenne : leviers de la gestion de crise ? » (12/2025)
\item Institut Montaigne -- Note d'enjeux -- Esprit de défense : l'affaire de tous (03/2025).
\item Lundi de l'IHEDN -- « L'IHEDN au cœur d'un écosystème stratégique unique » (03/2026).
\item Revue d'Anthropologie des connaissances -- Article « Techno-politiques des vulnérabilités et production des catastrophes et crises » (20/01/2026).
\item Hal open source -- Bulletin of Sociological Methodology  : « Un protocole pour une enquête par questionnaire anonyme au sens du Règlement européen ». Marie Plessz (2020).
\end{itemize}

\subsubsection{Rencontres d'acculturation à notion de crise et aux méthodes d'enquêtes}

\begin{itemize}
\item Risques/France 20230 - Programme de recherche : « Les SHS à l'épreuve des risques et des catastrophes ». Campus Saint-Germain (06/11/2025).
\item CNRS -- Projet RISC (Risques et sociétés à l'ère des changements globaux : enjeux, savoirs et politiques) : Quels renouvellement des études de catastrophes ? -- Participation et mobilisation citoyenne dans la gestion des risques et des catastrophes. IPGP Jussieu (17/12/2025).
\item CNRS -- Journées scientifiques annuelles PEPR Risques : Construire des sciences des risques intégrées : défis et difficultés -- Recherche, expertise et action publique : quelles nouvelles perspectives pour une réduction durable des risques de catastrophes ? » (26 -- 27 mars 2026)
\item CNRS /EHESS Condorcet : « Savoirs et gouvernement des situations de crises et de catastrophes, XIXe -- XXIe siècle : cadrage et instruments du gouvernement des catastrophes, mise en mémoire et oubli de l'histoire des catastrophes, les politiques de l'alerte et de la réduction des catastrophes (08/04/2026).
\item Paris I Panthéon La Sorbonne -- Journée d'étude - « Protéger l'enquête : enjeux scientifiques et éthiques ». Groupe de recherche AFSP. Victor Violier, Vanessa Frangville, Marielle Delbos (20/03/2026).
\item Forum de Paris pour la défense et la stratégie :  tables-rondes sur les thèmes : « Engagement majeur, nouvelle génération et «  ; « Le rôle de la société civile dans la défense et la préservation de la paix » ; « Préparer la nation au combat : les forces morales et la résilience de la nation » (24-26/03/2025).
\item Le Point -- Institut de France -- Forum « Guerres et Paix » : présentation de l'initiative « Au contact, citoyens, citoyennes » lancée par Muriel Domenach, ancienne ambassadrice de France auprès de l'OTAN (01/04/2026).
\end{itemize}

\section{Cheminement de la faisabilité de la démarche}

\subsection{Réassurance technique auprès d'universitaires}

\textbf{Rencontres avec:}

\begin{itemize}
    \item Victor VIOLIER, Sociologue IRSM (Pôle Défense et Société). Echange avec des membres du Comité le 20/02 et le 20/03 (Journée d'étude CNRS/EHESS/AFSP/Panthéon-Sorbonne -Protéger l'enquête : enjeux scientifiques et éthiques) :
    \begin{itemize}
        \item Points positifs
        \begin{itemize}
            \item L'enquête constitue un instrument de connaissance interne, un levier de structuration stratégique, un outil de valorisation institutionnelle, un fondement empirique pour toute proposition relative à l'intégration de l'AR 16 dans une situation de crise majeure
            \item Mise en perspective d'une séance de restitution associant L'AR 16 et L'IRSEM
        \end{itemize}
        \item Point de vigilance
        \begin{itemize}
            \item Clarté de la problématique, rigueur de la construction du questionnaire et respect du cadre juridique applicable
        \end{itemize}
    \end{itemize}
    \item Anne MUXEL, Directrice déléguée du CEVIPOF (Centre de recherches politiques de Sciences Po) et directrice de recherches émérite en sociologie et sciences politiques au CNRS, ancienne directrice du domaine « Défense et société » de l'IRSEM :
    \begin{itemize}
        \item Points positifs
        \begin{itemize}
            \item Démarche probante
            \item Structure et construction du questionnaire
            \item Le dispositif prévoit d'associer l'outre-mer
            \item Disponible pour accompagner la démarche
        \end{itemize}
        \item Points de vigilance
        \begin{itemize}
            \item RGPD
            \item Anonymisation du questionnaire à préserver
        \end{itemize}
    \end{itemize}
\end{itemize}


\subsection{b.	Réassurance interne par le biais d'échanges}

\textbf{Rencontres avec}

\begin{itemize}
    \item Catherine SARLANDIE DE LA ROBERTIE, Présidente de l'Union IHEDN, Préfète, Professeure des Universités, Rectrice d'Académie, Directrice de la Formation continue Panthéon-Sorbonne :
    \begin{itemize}
        \item Points positifs
        \begin{itemize}
            \item Démarche pertinente à valoriser dans le cadre des anniversaires en cours
            \item Dispositif à tester dans 3 régions dont une en outre-mer avant un élargissement à l'ensemble des AR
            \item Proposition à présenter l'étude réalisée à La Sorbonne (Paris 1) en septembre 2027
        \end{itemize}
        \item Points de vigilance
        \begin{itemize}
            \item Distinguer enquête et questionnaire
            \item Discussion possible sur l'anonymisation des données
        \end{itemize}
    \end{itemize}
    \item Pierre VUILLAUME, Général (2S), Délégué général de l'UNION-IHEDN :
    \begin{itemize}
        \item Points positifs
        \begin{itemize}
            \item Interet pour la démarche
            \item Proposition inédite
        \end{itemize}
        \item Points de vigilance
        \begin{itemize}
            \item Peu de ressources ou de références historiques pour documenter la préparation du projet
            \item Dispositif lourd
        \end{itemize}
    \end{itemize}
\end{itemize}

Les conclusions de tous ces échanges ont été présentés en comité et on fait l'objet de discussions.

\section{Cheminement de la décision politique de l'AR 16}

\begin{itemize}
    \item 6 mars 2026 : le comité d'études adresse une lettre-courriel proposant le projet au Bureau de l'AR 16 pour une mise en œuvre éventuelle du questionnaire destiné à révéler une cartographie précise de ressources et compétences des membres de l'association.

    \item 7 mars 2026 : la présidence de l'AR confirme la transmission du projet au Bureau avant la poursuite éventuelle de la discussion lors du CODIR du 17 mars.

    \item 19 mars 2026 : le comité d'études est informé de l'idée d'intérêt que suscite le projet de questionnaire au sein du CODIR. Un aval officieux du questionnaire est évoqué.

    \item 3 avril 2026 : un échange entre la présidente de l'AR 16 et la présidente du Comité d'études permet d'obtenir un accord pour le projet de questionnaire.

    \item 10 mai 2026 : la présidente du comité est informée de l'invalidation du projet de questionnaire.
\end{itemize}

\end{document}
